\chapter{C14.1b: the remaining source at bounded Alice-to-reciprocal-Bob ratio}
\label{component:p50b4df0bb619}
\noindent\textbf{Component }\path{p50b4df0bb619}\par
{\small\noindent Source: \path{BASE/proofs/c1_unbounded/c141b_gamma_angular_domination.tex}\par}
{\small\noindent SHA256: \path{50b4df0bb619d10389dbddf2ac9c3de366990af3f598368843141cec5fe02b81}\par}
{\small\noindent Inventory use: GA1-GA9 surviving bounded-u/w admission and exact remaining domain | GA1--GA9 | C14\_\allowbreak{}ENTIRE\_\allowbreak{}ORIGINAL\_\allowbreak{}PARENT\_\allowbreak{}GRAPH; precise subordinate uses retained below\par}
\medskip
\section{C14.1b: the remaining source at bounded Alice-to-reciprocal-Bob ratio}
\label{p50b4df0bb619:sec:c141b-gamma-angular}

\paragraph{Exact additional scope.}
Keep the full original C14.1b source, physicality, strict gaps,
nonnegative coefficients, total-scale and prior-complement
intersections. Use the independent fields and the complete P4/SP6
parent of the inherited \path{parent_reference.tex}. In this
section $h$ is a positive bounded unit and $a_0=k/\sigma$ is
infinitesimal. The unit-$a_0$ source was already admitted by
US1--US15; no conclusion about that proof is changed here.

\begin{theorem}
On the stated source, the entire original parent is tensor SOS
for every projective testing configuration whose remaining
double-small-angle chart has $u,w>0$ infinitesimal and $u/w$
bounded, and for every real or complex state. Include the exact
sign gauges, literal endpoints and other testing exits already
admitted in SF1--SF21. There is no restriction on a state ratio
or on $c=uW/(wZ)$ in this theorem.
\end{theorem}

This is an additional testing-domain admission. It does not close
the remaining chart $u/w$ unbounded. Unlike BC1--BC11, its testing
predicate is independent of the state, so the same predicate holds
for the real and imaginary state components.

\paragraph{Available source pivots.}
The source bounds SF1--SF5 and US1--US3 give
\[
 L_0=\lambda\sigma\sim2h,\qquad
 \lambda^{-1}=\sigma/L_0\text{ infinitesimal},\qquad
 \rho=L_0a_0^2\text{ infinitesimal},\qquad
 d=4D/\sigma^2\sim2.
 \tag{GA1}
\]
Indeed US2 first bounds $a_0$, and its exact residue identity makes
$b_0$ infinitesimal when $a_0$ is; the sphere then gives $d_0\sim1$.
Thus the original Gamma diagonal supplies a strict positive-unit
pivot. Also $\lambda x=\rho\sigma$ is infinitesimal. The source
coefficients $K_A,\chi,\nu,R_\gamma,R_x,A_y$ are bounded. No
lower cap on $\beta$, $b$ or $Q$ is imposed. A literal $D=0$ face
is incompatible with GA1, rather than being removed from a domain
where it could occur.

\paragraph{A paid angular domination from the original Gamma square.}
The unchanged sign and state-axis exits leave $u,w,W,Z>0$; their
ratios may otherwise be arbitrary. Set $t_s=\sqrt\sigma$ (a source
letter, unrelated to the earlier notation $t=u^2$), and define
\[
 E_0=quW-t_s Z,\quad V=uE_0,\quad
 \alpha_t=q\frac uw,\quad
 U=\frac{B_1-B_3}{\sqrt\lambda}
   =\frac{d_1-d_3}{\sqrt\lambda}.
 \tag{GA2}
\]
Here $B_i,G,d_i=B_i-G$ are exactly the four monomials and circuit
differences in US7--US10. Every root in GA2 is separated. The only
new testing quotient is $u/w$, which is bounded by the theorem's
predicate. Hence $\alpha_t$ is a bounded infinitesimal coefficient.
Exactly
\[
 \Gamma+t_sR=-\frac{w+t_sq}{q}E_0,
 \qquad A_\theta=U+V.
 \tag{GA3}
\]
The first identity is used as a polynomial identity; its mixed
factor $w+t_sq$ is never inverted. In fact it yields the explicit
separated-SOS difference
\[
\begin{split}
 \Delta&:=\alpha_t^2(\Gamma+t_sR)^2-V^2\\
 &=u^2E_0^2\left\{\frac{2t_sq}{w}
                            +\left(\frac{t_sq}{w}\right)^2\right\}\\
 &=\left(\sqrt2\sqrt{t_s}\sqrt q\,
                  \frac{u}{\sqrt w}E_0\right)^2
   +\left(\frac{t_squ}{w}E_0\right)^2.
\end{split}\tag{GA4}
\]
All square roots and inverses in the last line belong to their
individual fields. The roots themselves need not be bounded.

Choose a positive real $C_*$ so that the bounded pure-source
coefficient $K_A+C_*$ has positive real residue. The following
once-paid identity allocates the complete signed angular diagonal:
\[
\boxed{\begin{aligned}
 K_AA_\theta^2={}&(K_A+C_*)A_\theta^2
  -\frac{2C_*}{\lambda}(d_1-d_3)^2
  -2C_*\alpha_t^2(\Gamma+t_sR)^2\\
 &+C_*(U-V)^2+2C_*\Delta.
\end{aligned}}\tag{GA5}
\]
It follows simply from $A_\theta=U+V$ and GA4. The last two terms
are retained SOS. The two negative terms have bounded
infinitesimal coefficients on the circuit and $(\Gamma,R)$ frames.
This constructs a strict angular pivot without requiring a positive
residue of $K_A$, and without the collapsed $\lambda k$ pivots of
US11.

\paragraph{Complete scalar, every state normal, and all excesses.}
Use the whole four-monomial allocation once:
\[
       \mathscr S=\sum_{i=1}^4d_i^2+\mathcal E_4
                                           -\lambda xR^2.
       \tag{GA6}
\]
Here $\mathcal E_4$ is the explicit separated SOS in US10. Do not
also allocate US6 or US11. The single radial cost in GA6 is an
infinitesimal pure-source coefficient.

Use SF8/CT7--CT7a with the full source-opening-safe radial squares
and every positive excess restored. The resulting strict bounded
radial frame is $(N_y,R)$, uniformly through the $b,Q$ openings.
Together with the original $N_x$ pivot, $d\Gamma^2$ from GA1,
GA5 and GA6, the principal frame is
\[
       (N_x,N_y,R,\Gamma,A_\theta,d_1,d_2,d_3,d_4).
       \tag{GA7}
\]
Its residue is strictly positive definite: the $N_x$ residue is
$2$, the Gamma residue is $2$, the angular coefficient is
$K_A+C_*$ with positive residue, and the four circuit coefficients
are $1$. All costs in GA5--GA6 and all CT7a remainders are bounded
infinitesimal matrix entries on this frame.

The two complete Gamma normals remain exactly
\[
      -4\nu N_x(u\Gamma)
       +R_\gamma(N_y+bR)(qw\Gamma).
       \tag{GA8}
\]
Their coefficients are infinitesimal directly because $u,q,w$ are
infinitesimal and the remaining factors are bounded. No division
by $\sqrt\rho$ is used. The five remaining non-Gamma crosses are
precisely SF9, with bounded infinitesimal coefficients on GA7.
The whole $N_yR$ coefficient has already been included in the
exact radial extraction. Thus every original SP6 term, including
both first state normals, is accounted for.

Keep separately and exactly once $4\beta D_wN_x^2$,
$2(2D+1)(uX_e)^2$, the nonconstant positive $N_x$ terms,
$4bpzu^2\Gamma^2$, all radial and source-weight excesses,
$\mathcal E_4$, and the last two SOS terms of GA5. The separated
bounded-Gram lemma applied to GA7 proves the complete parent.

The exact inherited literal, sign and remote exits supply the
complementary charts specified in the theorem. For complex states,
apply the same real theorem separately to each real state part;
the hypothesis $u/w$ bounded depends only on the test and is
unchanged. Restore only positive pure-source and pure-testing
normalizations. No source/testing rate is identified or compared.

\paragraph{Exact updated residual and failed extension.}
Combining this theorem with SF, BC and US leaves C14.1b only on
the original retained complement intersected with
\[
\begin{gathered}
 \operatorname{res}h\ge\tfrac12,\quad k/\sigma\text{ infinitesimal},
 \quad b_0\text{ infinitesimal},\quad d_0\sim1,\\
 u,w>0\text{ infinitesimal},\quad WZ>0,\quad W/Z\text{ unbounded},\\
 c=uW/(wZ)\text{ unbounded},\qquad u/w\text{ unbounded}.
\end{gathered}\tag{GA9}
\]
The exact obstruction to this particular extension is
$\alpha_t=q(u/w)$ in GA2: its testing factor is no longer bounded.
It cannot be declared small by comparing the independent source
$q$ to the testing reciprocal $w/u$. At contact
$u_c/w_c=a_0\sqrt\sigma/q$; a further lawful source split by
$q/(a_0\sqrt\sigma)$ may be useful, but no all-tests conductor or
join for that split has been proved here. GA9 is not asserted
nonempty, and full C14.1b remains open.

The focused checker verifies GA3--GA5 and the complete scalar
budget in GA6. Analytical source bounds, every parent term and the
exact quantified testing join are the proof above; this remains
same-system mathematical review.
